\daymark{2}{Stone 8, Y.H. 812 (Eighthday)}

\dropcap{T}{he kiln} had been breathing since the middle of the night. In the \place{Vail kiln-house} at first light a man waking in the next room would have known the place by that sound and by nothing else: a low intake drawn through a throat of brick, a pause as though the thing had forgotten what it was about, and then a letting-out so long it went on past the point where a breath had any business ending. It was the breath of the fire. It had been kept by coals banked under ash for two years, and it would go on after the house had stopped speaking and gone to its work. People who lived beside it had long since ceased to hear it, the way a man ceases to hear his own breathing until it fails.

It came in fours. Four draws, four draws, and then a space where the fifth should have been, a gap the width of a thumb in the order of the thing, and then the long letting-out that filled the room with smoke and went on until you thought the coals had given all they had and found another breath in them. \person{Tarn}, who had slept through the river icing over and the house growing cold around him, had never counted it. \person{Sella Vail}, who had kept the fire through two winters without missing a night, had never needed to. It was simply how the house breathed, and a man might live beside it all his days without hearing the missing fifth, and then one morning the missing fifth would be all he could hear.

You could have known the house without ever seeing a face in it. In the \place{Vail wheel-shed} the wheel stood with its seat worn hollow on the left, a shallow bowl in the wood where a man had sat for twenty years and leaned left into the throw, and the hollow was empty now and would stay empty; Sella had twice stopped a hand that went to fill it. Behind the door a coat on a peg kept the shape of a broad man, salt stiff in the seams from a winter of carrying slip, and the shape was broader than any man who lived in the house now. \person{Jorra Vail} had reached for it in the dark twice and stopped both times and said nothing about either. Along the skirting board ran a ridge of \lore{the Ash}, the width of a thumb and thicker than the same ridge in any other house in \place{Low Kiln}, sour in the nose when the room was warm, plain as a line struck in chalk when it was cold. Nobody swept it. It would be swept on \lore{the Last Day}, with the rest of the district's sweeping, and burned next morning and called dealt with, which was how Low Kiln dealt with most things it did not want to look at.

When the kiln was low and the house had gone quiet, the smell of him was stronger than any memory of him. It came out of the corner where the Ash lay and out of the coat on its peg and out of the hollow in the seat: wet clay and banked coals and the sweetish edge of lead glaze, gone thin and cold, the way a room smells where something has been left too long and nobody has moved it. Jorra had been thirteen when his father went and had taken the wheel a week later, and had spent three years trying to fill the hollow with his own weight, and had not once said that he could not. The wheel did not care whose hands worked it. It cared if they were steady.

The sayings of \person{Garl Vail} had outlived him and improved in the keeping, as sayings do when nobody remembers the tone they were said in. The house kept one of them and said it over the pot and the fire and on mornings when none of the three people in it had slept. “Leave the crack showing. People trust a thing they can see has been handled.” He had been a big man with clay in the creases of his hands and a laugh that ended in a cough, and had gone in the Long Cold when there was no crossing the river for a collector, and had not been collected. His warmth lay where it had fallen. The ridge along the skirting board was the shape of it; the coat kept his shape on the peg; the seat kept the leftward lean of him in the wood. None of these facts had been entered in any account a man come down from \place{Slate Row} with a slate in his hand could read.

The house kept a roll of its own, and had kept it since Garl was a boy, and did not show it to that man either. On its private slate Garl was entered and Enna was entered, and the foundling's name, when he came, would be entered in a hand that was careful about such things. It was a small thing kept close, and whether it was a true roll or a comfort nobody in the district had ever asked, and the house had never offered to say. A household that keeps its own account does not need to explain the keeping; it needs only to keep it.

Tarn had been sweeping since before anyone woke. It was not in the order of the house for a clay-bound hand to be up before the fire was stirred, but Tarn had never needed ordering in that particular, having learned early that a broom in his hand was a thing nobody could call him for. Six years of a seven-year term lay behind him and he had counted every one, and would not, he said, count the next. He called Sella “house,” the way he called everyone house, and said “I owe” where another man said “I am.” He swept the clay dust toward the door and stopped. A pot had been set down somewhere behind him and the small sound of it went through the room and into his sweeping, and for the space of two strokes the broom did not move. Then it moved again. He had not turned his head. Nobody in the room had said anything about the pot.

When he thought nobody was looking he put the small bench knife into his pocket. It was a stub of a thing with the mark on its handle worn almost flat from use, and it went in quickly and without ceremony and was gone. Then he said, of his own accord and to nobody in particular, that he had taken the small knife and put it back and they could check if they liked. He said it in the tone of a man stating the weather, and the tone invited a laugh that did not come. His deference had a nail in it: he agreed so thoroughly that an argument had nowhere to go, and the thoroughness of it was how you knew where he stood. Sella had counted the tools before she opened her eyes that morning and did not look up. The house needed no telling.

Jorra stood at the wheel in the shed with the door rolled back and had been there since the coals were stirred. He was sixteen and built like a hinge, tall for a clay boy and narrow, with slip dried to the colour of old bone up both forearms and a permanent bend in the left little finger. He told people the finger came from a fight, because a fight sounded more like him; it had come from the wheel when he was eleven. He did not look at Sella when she spoke to him. He never had. He obeyed her instantly, which was the other half of the same habit: a son who cannot bear to look at a command and cannot think of disobeying one. He spoke to the wheel first and to the pots after it and to people last, and to people in short flat declaratives that stopped where a longer sentence would have had to go on.

“Is it dry?” he asked the pot. To the wheel he gave the little click of the foot that started it, and to Sella, when the work required it, he said, “The slip's too thick. It'll crack at the shoulder.” He had a way of going quiet when \lore{the levy} was spoken of that the household had learned to read as weather, and there was a temper under the quiet that nobody had yet seen the bottom of. Men had come into the back room through the winter and spoken low about slates and accounts, and Jorra had listened to all of them and given none of them a word they could carry away. When they pressed him he said he had not said he was with them and had not said he was against them, and that in this district those were the same thing and he had said one of them anyway. He would not say which.

Sella kept the house the way a widow keeps a wound clean: daily, without comment, with a hand that knew exactly how much the thing would bear. She was not tall and did not make herself tall; she took up room by being impossible to move. Before she opened the kiln she put her thumb to the cord at her throat and to the \lore{widow's mark} hung there, a thumb-length of lead-glazed clay with the Vail kiln scratched into it, and held it for the count of a breath, and let it go. She touched it at no other time. It was not a thing she did when she was afraid; she did it before the fire, and there were people in the district who said the fire would go out if she ever forgot. She would have let them think it. When she was angry she got quieter, and the quieter she got the more the house straightened itself up around her, as it was straightening now.

“Feed it,” she said to the fire, and the fire was fed. She said “the house” where another woman said “I,” and gave instruction to the clay and the coals and never a finished feeling, because a thing said aloud was a thing the Court could enter, and she had learned to keep her feelings where no Court's hand could reach them. She would not let anyone fill Garl's seat or take his coat down from the peg, and had never given a reason for either, and the household had stopped asking for one years ago.

There was a foundling in the district with no house-name, three doors down the yard in \place{House Silt}, and a slate in that room with a name in a hand nobody there could read. The boy was to be brought under this roof on the morrow, the third day, and nobody in the kiln-house had said what would be paid for him or why, and nobody had said his name. It had been decided without a meeting, the way such things were decided in Low Kiln: a coat left on a peg in a poor house is a coat anybody may wear, and who left it is a question nobody is obliged to answer. A stone warmed at this fire, the colour of skimmed milk, had already crossed the yard in an apron and lay cold now under a fold of wool in the room three doors down, and a crust of brown-edged barley carrying the smoke of this kiln had gone the same way. Neither had been entered against anybody. The house that had paid for both had not once said so.

Sella stood with her hand on the kiln door and said the thing that had been waiting in the room all morning. “The boy comes under the roof on the morrow,” she said. “The house will keep him.”

Jorra did not look at her. His hands went on turning the pot, the left little finger bent and steady, and the wheel went on with its low small sound under his palms.

“He's not Vail,” he said.

It was not loud. It was not an argument in the way arguments are usually made, with the voice up and the face turned; it was the flat declarative of a boy who had put a thing where it could not be taken back, and who knew that in this house a thing said was a thing entered. Sella did not answer it. She had been quiet before and she got quieter now, and the quiet went into the kiln-house like cold into a room, and Tarn's broom found the dust and swept it, and the wheel went on turning.

In the ash where the fire had been cleaned the night before, the charred edge of a slate lay among the coals, pushed in at some time and never raked out, the writing gone and the edge gone black. Nobody had picked it up to look. The house had swept around it for a year without once asking what it was, and had gone on keeping its own account.

She took her hand from her throat and put it on the kiln door and opened it. Heat came up out of the fire and took a degree of chill from the room, and woodsmoke went into the wet wool of the coats and stayed. “Don't set that near the door,” she said to the pot. “You'll pay for it in the thaw.” She gave the fire what it needed and said nothing further, and the house went on.

Three doors down the yard, in a room with a door that would not shut, a child lay under a fold of grey wool with a spent stone at his belly and did not know he had been argued about. He had not yet come under this roof. The third day was still a day off. \place[the Ilm]{The Ilm} lay under its skin of ice, and \lore{the Five} were spoken of as three days off, and Low Kiln said three days when it had counted what it had and did not want to say it. In the kiln-house the wheel went on turning and Tarn swept and Jorra did not look up, and the words stood where he had put them, in the space between a mother and her son where the foundling would lie when he came.
